\section{\tool{} Design}\label{sec:design}

Section~\ref{sec:study} shows that binding flaws are decided by comparing
two artifacts that no single analyzer reads together, and that auditors find
them by reasoning about intent. \tool{} mechanizes that reasoning. An LLM
\emph{brain} reads the target. A fixed \emph{skill procedure} forces the brain
to build the two ledgers of Section~\ref{sec:def-ledgers} explicitly, row by
row. A set of \emph{analysis tools} supplies facts the brain should not guess.
\DEMO{An \emph{exploration loop} then lets the brain work like an auditor writing a
proof of concept: it plans an attack on each warning, writes it as a test,
runs it in a sandbox, reads the outcome, and tries another path when the test
fails. Each warning leaves \tool{} as an \emph{evidence package} that a human
can check in minutes.} Figure~\ref{fig:arch} shows the architecture.

\begin{figure*}[t]
\centering
\resizebox{\textwidth}{!}{%
\begin{tikzpicture}[font=\scriptsize, node distance=0.35cm and 0.5cm,
  st/.style={draw,rounded corners,align=center,minimum height=0.85cm,text width=2.05cm,fill=gray!8},
  tl/.style={draw,align=center,minimum height=0.6cm,text width=1.75cm,fill=blue!6},
  ev/.style={draw,align=center,minimum height=0.85cm,text width=1.8cm,fill=green!8},
  ar/.style={-{Latex}}]
\node[draw,align=center,text width=1.6cm,minimum height=1cm] (t) {Target\\(circuit, verifier, consumer)};
\node[st,right=of t] (s0) {0. Inventory\\entries $E_i$, statements $S_j$};
\node[st,right=of s0] (s1) {1. Effect ledger\\(entry, effect, value, verdict)};
\node[st,right=of s1] (s2) {2. Statement ledger\\(public value, meaning, constraint)};
\node[st,right=of s2] (s3) {3. Pattern checks\\P1--P5 playbooks, follow-through};
\node[st,right=of s3] (s4) {4. Attacker pass\\and final review};
\draw[ar] (t)--(s0); \draw[ar] (s0)--(s1); \draw[ar] (s1)--(s2); \draw[ar] (s2)--(s3); \draw[ar] (s3)--(s4);
\draw[ar,dashed] (s2.south) to[bend left=25] node[below,font=\tiny]{coverage re-prompt} (s1.south);
\node[tl,below=0.75cm of s0] (a) {inventory: entries, effects, publics};
\node[tl,below=0.75cm of s1] (b) {verifiers, pins, decode};
\node[tl,below=0.75cm of s2] (c) {def / callers / callees};
\node[tl,below=0.75cm of s3] (d) {replay, P4 checks};
\node[draw,dashed,fit=(a)(b)(c)(d),inner sep=3pt,label={[font=\tiny]below:analysis tools (deterministic)}] (tools) {};
\foreach \x in {a,b,c,d} \draw[{Latex}-{Latex},blue!50] (\x) -- (\x |- s0.south);
\node[ev,right=0.6cm of s4,yshift=0.9cm] (e1) {Warning\\(ledger row, reasoning)};
\node[ev,below=0.25cm of e1,dashed,draw=red] (e2) {Exploration loop\\plan, test, run, retry};
\node[ev,below=0.25cm of e2,dashed,draw=red] (e4) {Evidence package\\(paths tried, executed test)};
\draw[ar] (s4) -- (e1.west); \draw[ar] (e1)--(e2); \draw[ar] (e2)--(e4);
\end{tikzpicture}}
\caption{\tool{} architecture. A skill procedure drives the LLM through
explicit ledgers (top). Deterministic analysis tools answer factual queries
(bottom). For each warning, an exploration loop writes attack tests, runs
them in a sandbox (local chain, the project's own prover), and attaches the
paths tried and any executed test. Humans adjudicate the package; no LLM
verdict serves as ground truth. \DEMO{Dashed: automation in progress.}}
\label{fig:arch}
\end{figure*}

\subsection{Why a procedure and not a prompt}\label{sec:design-why}

A plain LLM asked to ``find binding bugs'' reads the consumer, notices the
most salient issue, and stops. On our test set it finds 3 of 14 known flaws
(Section~\ref{sec:eval}). The failure is not knowledge but coverage: a target
has many entry points, each with several effect values, and a flaw hides in
one row of a table that the model never writes down. The skill procedure turns
the definition into that table and makes omissions visible
(Figure~\ref{fig:controller}). Every step emits
structured rows (\texttt{BINDCHECK\_LEDGER}, \texttt{BINDCHECK\_STMT}), and the
controller checks the rows against the inventory before it lets the brain
proceed.

\begin{figure}[t]
\begin{lstlisting}[mathescape]
Scan(target):
  copy target read-only; I := Inventory(copy)   // E_i, S_j
  L_E := Brain(step1, I)                         // effect ledger
  for r in 1..2: U := I.E \ entries(L_E)
      if U = {}: break;  L_E += Brain(reprompt, U)
  L_S := Brain(step2, I, L_E)                    // statement ledger
  for r in 1..2: U := I.S \ values(L_S)
      if U = {}: break;  L_S += Brain(reprompt, U)
  W := Brain(step3, L_E, L_S)  // playbooks P1-P5, follow-through
  if cost < SOFT and time < T_soft: W += Brain(step3b, W)
  W := Brain(final, W)         // drop C1-C3 failures, 1 per root cause
  for w in W:                     // exploration loop
    for k in 1..K: t := Brain(attack, w, trace(w))
      o := Sandbox(t); trace(w) += (t, o)
      if Effect(o.attack) != Effect(o.honest): break
  return [(w, trace(w)) for w in W]   // evidence package
\end{lstlisting}
\caption{The \tool{} controller. Every \texttt{Brain} step returns
structured rows that the controller checks against the inventory
before it continues. The sandbox, not the brain, computes whether an attack
test changed the effect.}
\label{fig:controller}
\end{figure}

\subsection{Skill procedure}\label{sec:design-skills}

\para{Step 0: Inventory}
A deterministic prepass lists the entry points that commit effects ($E_1,
E_2,\dots$) and the producers of public statements ($S_1, S_2,\dots$): the
verifier calls, the public-input arrays, and the circuit or guest-program
outputs. The list bounds the work and becomes the coverage target.

\para{Step 1: Effect ledger}
For each entry $E_i$ the brain writes one row per (entry, effect, value). The
row names the effect (transfer, mint, state write, grant), the value that
decides it, where the value comes from (public input, hash-committed input,
storage, calldata, \texttt{msg.sender}), and a verdict: \emph{bound},
\emph{unbound}, or \emph{needs statement}. If an entry has no row, the
controller re-prompts for it, at most twice.

\para{Step 2: Statement ledger}
For each public value the brain records what the circuit or guest constrains
it to mean, which consumer value it is supposed to fix, and whether the
consumer actually compares the two. This step catches what the consumer alone
cannot show. A recipient that the circuit hashes into a commitment is bound
only if the consumer recomputes and checks the same commitment.

\para{Step 3: Pattern checks and follow-through}
For every unbound or unresolved row, the brain runs the playbook of the
matching pattern: recipient and fee provenance (P1); chain, contract, root and
prior-state checks (P2); replay-key derivation and field aliasing (P3);
verifier result handling, key and program-ID pinning (P4); public-input
order, width and encoding (P5). It then follows the value into libraries,
helpers and record writers. Many flaws in our dataset sit one call below the
entry point.

\para{Step 4: Attacker pass and final review}
The brain takes the view of an attacker who holds an observed valid proof and
asks what it can change. The final review drops warnings that fail C1--C3,
for example admin-only paths or self-inflicted denial of service, and emits
one warning per root cause. Each warning carries the pattern, the consumer
file and function, the effect, the unbound value, an explanation, and the
related locations (verifier, circuit, helpers).

\para{Where the skills come from}
The procedure and the five playbooks were written from the development split
only (Section~\ref{sec:eval-setup}). We split the development families again,
with a fixed seed, into a \emph{tune} part and a \emph{check} part. In the second
design round, changes were motivated by misses on the tune part. A change was
kept only if it did not hurt the check part. The check part is not a clean
holdout. The first version's skills were written on all development items,
and two check-part misses of that version were named in the second round's
brief. On development data, the frozen version finds 9 of 13 flaws at the
exact location and mechanism, 5 of 6 on the check part, up from 4 of 6 for
the first version. The held-out evidence is the test set
(Section~\ref{sec:eval}).
% src: mvp/tool/FREEZE-v2.md (ALL 9/13; CHECK 5/6; v1 CHECK 4/6)

\subsection{Analysis tools}\label{sec:design-tools}

The brain can query deterministic tools instead of guessing. These tools
include:
\begin{itemize}
\item an inventory of entry points, effects and public-input arrays;
\item a verifier locator that reports how each verifier is called and whether
its result is checked;
\item a key and program-ID pin check;
\item a public-input decoder;
\item a replay-key extractor;
\item cross-file definition, caller and callee lookup.
\end{itemize}
The current tools work at the syntactic level, over contract and circuit
sources alike.
\DEMO{An AST-based extractor replaces the syntactic inventory for Solidity and
Circom and resolves public-input indices across the verifier call and the
circuit's \texttt{main} signature.}
\TODO{A1}{AST-based extractor for Solidity/Circom public-input index resolution}

\subsection{Exploration loop and evidence package}\label{sec:design-explore}
\DEMO{A warning on its own is a claim. Auditors do not stop at a claim: they write a
proof of concept, and when one attack path fails they try another. \tool{}
gives the brain the same loop. For every warning that survives the final
review, the brain plans an attack, writes it as a test, and passes it to a
sandbox, which runs it and returns the outcome. The brain reads the outcome
and either stops or revises the attack, for up to $K$ attempts. Failed paths
are kept, because they tell the auditor what was ruled out.}

\para{The sandbox}
\DEMO{Every attack test submits the same accepted proof twice on the same state:
once as the honest call, and once from a second account with the
attacker-chosen value of the suspect effect row. The test succeeds when both
calls are accepted and the committed effects differ in recipient, amount,
state written, or count. This is condition C3 of
Section~\ref{sec:def-ledgers}, made executable. The sandbox offers two
strengths of proof. A \emph{mock-verifier test} replaces the verifier with one
that accepts exactly the honest public-input vector. A \emph{real-proof test}
generates the honest proof with the project's own prover and submits it to the
unmodified verifier and consumer on a local Anvil chain. No test touches a
live network.}

\para{Evidence package}
\DEMO{\tool{} reports every warning with (i) the ledger row and code location,
(ii) the brain's reasoning, (iii) the trace of attack paths tried, failed
ones included, and (iv) the executed test and its output, when one succeeded.
We label a package by its strongest content: \emph{E1} (claim only),
\emph{E2} (a mock-verifier test succeeded), or \emph{E3} (a real-proof test
succeeded).}

\para{Who decides}
\DEMO{The brain may judge its own results, and \tool{} uses that judgment to drop
and rank warnings. No number in this paper rests on it. Whether a test
succeeded is computed by the sandbox from the two executions, not read from
the brain's summary, so a test that asserts nothing, or one that mocks away
the check under test, cannot succeed. Ground truth comes only from the audit
labels and from human adjudication of the evidence package
(Section~\ref{sec:eval-setup}).
The exploration loop and the sandbox support Solidity consumers with
Circom, gnark and Noir circuits.}
\TODO{A2}{fold the frozen validator v1.1 drivers (mock verifier, per-toolchain proof drivers) into the brain's toolbox as a sandbox tool; add the exploration loop with K attempts}

\subsection{Guards and isolation}\label{sec:design-guards}

Each scan runs on a fresh, read-only copy of the target. The brain gets only
read-only tools and has a 20-minute limit. The controller tracks the cost
of the session. Past a soft limit it skips the optional passes, and past a
hard limit it asks for the final answer at once. The brain is a commercial
model (Claude Sonnet, high reasoning effort). Any model that supports tool use
can be swapped in.
\TODO{E7}{optional: open-weight brain on B1/B2 (lab machine)}
